% =====================================================================
%  Chapter 24 — The Proof in Natural Language, 8
% =====================================================================

\chapter{Automorphy, the Tower, and the Vanishing}
\label{ch:assembly}

The last chapter of the proof assembles the pieces and states the two places
where the argument does not close.

\section{Automorphy of the characters}The characters $\mu_j$ that index the stages of the chain must themselves be
automorphic, and this is where the remark's most opaque device appears.

\begin{quote}
\small
``$\mu$ is automorphic from $\theta$ on sixfold cover of $\mathrm{GL}_5$ as a
period integral of thetas, quadric has natural interpretation on corresponding
CM e.g.\ $\mathbb{C}/\mathbb{Z}[\omega]$ as Weil, put Poincar\'e bundle on the
edges to get $\theta$s.''
\end{quote}

\begin{quote}
\textbf{Proposition AB.} The characters $\mu_j$ are the Fourier coefficients
of automorphic forms, obtained by taking the theta series of a line bundle on a
sixfold cover of $\mathrm{GL}_5$ in which the quadric has its natural
interpretation on the corresponding CM curve.
\end{quote}

The formal counterpart is the subtree
\code{Development/AutomorphicArithmetic/AutomorphicVarieties/TheoremsMuLinked/},
and beneath it the deepest input of the whole development,
\code{GlobalTheory.Mem.LiteralAutomorphy.LitN} --- the ``literal
automorphy'' that the bundle \code{Af.Hypotheses\_2} takes as its first field.

\section{The relative trace formula}The trace formula supplies the transfer of control from the group where the
analysis is easy to the group where it is needed.

\begin{quote}
\textbf{Proposition AC.} The $\mathbb{Z}_p$-Galois spectral sequence for an
inner form $G$ of $G'$ gives
\[
H^*(\mathbb{G},\mathbb{F}_p) \Longrightarrow
H^*_{f\text{-}is}(G,\mathbb{F}_p) \oplus H^*_{f\text{-}c}(G,\mathbb{F}_p),
\]
and the arithmetic of the Selmer groups of $G$ is controlled by the continuous
cohomology $H^*(G',\mathbb{F}_p)^{\mathcal{H}}$.
\end{quote}

The remark's phrase is that the trace formula is being used ``in a different
half-plane argument'', which means: not in the usual direction.

\section{The Poincar\'e series and Cauchy--Schwarz}\begin{quote}
\textbf{Proposition AD.} The size of the space of lifts is controlled by the
norm of the Poincar\'e series of the unitary group; the norm is estimated by
Cauchy--Schwarz, and the second factor is controlled by the approximate
orthogonality of the columns of the integral structure.
\end{quote}

The approximate orthogonality is what the averaging of the previous chapter
earns: the columns are the characters, and the equidistribution of
\code{TranslateMean} is what makes them near-orthogonal.

\section{The tower, and the three compatibility conditions}\begin{quote}
\textbf{Proposition AE.} The volumes of the quotients, the index of the theta
series at each stage, and the congruences relating the index at one stage to
the next, are mutually consistent; and the Iwasawa tower supplies the
arithmetic growth at every level.
\end{quote}

The three items are the remark's item 18, and the tower is item 19.

\section{The vanishing step}And now the obstruction.

\begin{quote}
\textbf{Proposition AF.} The Selmer group of an admissible representation
vanishes: for every odd prime $\ell$, every CM field $\Omega$, and every
admissible rank-two representation $W$ over the $\ell$-th cyclotomic
extension, the Bloch--Kato Selmer group of $W$ is trivial.
\end{quote}

This is \code{VanishOrd}, with the seven-field admissibility condition
\code{admitted}: continuity, determinant $\pm 1$, oddness under conjugation,
irreducibility, level away from $\ell$, the modular condition with a newform
and a character satisfying \code{KillsMu} and \code{Newform.Att}, and the flag
condition at a place above $\ell$.

The mechanism is twisting. Selmer groups are killed by cup product with a
character; the part that is not killed is the $H^1_o$ part, which is the
obstruction-theoretic part of $H^1$. The remark's item 17 says that what
survived the twisting was $H^1_o(\mathbb{Q}_\ell/\mathbb{Z}_\ell(-e))$, and
that killing it would require ``a significant weakening of Vandiver''.

\begin{quote}
\textbf{Vandiver's conjecture} asserts that the plus part of the class group of
$\mathbb{Q}(\zeta_m)^+$ has odd order for every $m$. It is verified for all
$m$ with $\varphi(m) < 2^{31}$ (Hart, Harvey and Ong, \emph{Mathematics of
Computation} \textbf{86} (2017), 3031--3049) and it is open.
\end{quote}

\section{Assembling the proof}\begin{quote}
\textbf{Proposition AG (the reduction).} Suppose that:
\begin{enumerate}
\item the determinant of Chapter~\ref{ch:determinant} is a Gram matrix with
the required interlacing sign structure and is non-degenerate;
\item the excess $R$ of Chapter~\ref{ch:count} has square-root quality;
\item the budget condition of Chapter~\ref{ch:budget2} holds with the factor
$2/3$;
\item the Selmer groups of Chapter~\ref{ch:assembly} vanish.
Then every non-trivial zero of $\zeta$ lies on the critical line.
\end{enumerate}

The chain of implication is: (2) gives the Weil bound's error term; (1) gives
the positivity of the Weil expression; Weil's criterion gives RH; the functional
equation and the classification of the trivial zeros give the statement about
the zeros of $\zeta$ on the critical line.
\end{quote}

\section{The two gaps, stated together}\begin{quote}
\textbf{Proposition AH (where the proof stops).} The argument requires, at
minimum:
\begin{itemize}
\item the factor $2/3$ in the budget condition to be replaced by $1$, or an
equivalent additional third of the budget to be found;
\item the $H^1_o$ classes to be killed without a full proof of Vandiver's
conjecture, or a weakening of it strong enough for the purpose to be
established.
\end{itemize}
Neither is supplied by the remark, and neither has been supplied by anything in
the development that has been audited.
\end{quote}

\begin{seamnote}{Where the formalisation ends}
The last formal link in the chain is \code{UpperTheory.Gal.Int.
galoisLocalGlobalStrong\_recParameter}, which is proved from
\code{galoisRepresentations\_of\_thetaLift}, which is proved from
\code{galoisLocalGlobalStrong\_of\_geometry\_lift}, which is proved from the
geometric realisation of weight one, which is proved from the Lefschetz fixed
point formula for a finite étale cover. Chapter~\ref{ch:chain} traces that
path in full.
\end{seamnote}

\begin{keytakeaways}
\item The characters indexing the chain must be automorphic, obtained from
theta series on a sixfold cover of $\mathrm{GL}_5$; this is the deepest
input in the development.
\item The Selmer groups must vanish; the formalisation states this precisely
as \code{VanishOrd} with a seven-field admissibility condition.
\item The two things the proof needs and does not have are the third of the
budget and the killing of the $H^1_o$ classes.
\end{keytakeaways}
